\documentclass[11pt]{article}
\usepackage{amsmath,amssymb,amsthm}
\newtheorem{theorem}{Theorem}
\newtheorem{lemma}{Lemma}
\title{Underloaded exterior interpolation with equal positive atom weights}
\author{}\date{}
\begin{document}\maketitle

\paragraph{Status.}
The following is a complete written local interpolation theorem. Its
rational identities, interpolation systems, and root-preserving endpoint
certificates have exact finite checks. It is not Lean-checked and does not
supply a recursive compiler. In particular it makes no assertion that
independently prescribed descendant patterns survive the split.

Let $P(t)=\prod_{h=1}^M(t-p_h)$, where the real poles are distinct and
belong to an open interval $(a,b)$. Let $1\leq q\leq M$, and let
$b<t_1<\cdots<t_q$ be exterior query times.

\begin{lemma}[Derivative interpolation]
The map
\[
 \{R\in\mathbb R[t]:\deg R<q\}\longrightarrow\mathbb R^q,
 \qquad R\longmapsto\bigl((R/P)'(t_j)\bigr)_{j=1}^q
                                                               \tag{1}
\]
is a linear isomorphism.
\end{lemma}
\begin{proof}
Put $f_k=(t^k/P)'$, for $0\leq k<q$. Write $W_m$ for the Wronskian
of $f_0,\ldots,f_{m-1}$, with derivative orders indexing rows.
The common-factor rule for Wronskians gives
\[
 \begin{aligned}
 W_m&=W(1,1/P,t/P,\ldots,t^{m-1}/P)\\
 &=P^{-(m+1)}W(P,1,t,\ldots,t^{m-1})\\
 &=(-1)^m\left(\prod_{k=0}^{m-1}k!\right)
                    \frac{P^{(m)}}{P^{m+1}}.                 \tag{2}
 \end{aligned}
\]
The first identity expands the Wronskian along its constant first
column. The common-factor rule follows by the triangular Leibniz rule
on derivative rows. The last identity expands along the final row:
all its entries except $P^{(m)}$ vanish, and the polynomial derivative
block has diagonal $0!,1!,\ldots,(m-1)!$.

For $t>b$ every factor of $P$ is positive, and the product rule gives
$P^{(m)}(t)>0$ for $1\leq m\leq M$. Thus every initial Wronskian
in (2) is nonzero throughout the exterior interval.

Here is the elementary zero-count implication used in this argument.
For $d$ smooth functions with nonvanishing initial Wronskians on an
interval, every nonzero linear combination has at most $d-1$ distinct
zeros. Induct on $d$. The first function $f_0$ never vanishes. Divide
a proposed combination by $f_0$ and differentiate. The new $d-1$
functions $(f_i/f_0)'$, $1\leq i<d$, again have nonvanishing initial
Wronskians, since
\[
 W\bigl((f_1/f_0)',\ldots,(f_m/f_0)'\bigr)
                  =W(f_0,\ldots,f_m)/f_0^{m+1}.
\]
If the derivative is a nonzero combination, Rolle's theorem and
induction give the assertion. If it is zero, the original combination
is a constant multiple of $f_0$, and has no zero unless identically
zero. Nonzero full Wronskian excludes an identically zero nontrivial
combination. Applying this implication to (2), a nonzero $R$ of
degree less than $q$ cannot give $q$ zeros in (1). The square linear
map (1) is therefore invertible.
\end{proof}

\begin{theorem}[Equal-weight underloaded split]
For every $\sigma\in\{-1,1\}^q$ there is a polynomial $R$ of degree
less than $q$ such that, for every sufficiently small $\delta>0$, the
two polynomials
\[
                       P_+=P+\delta R,\qquad P_-=P-\delta R
                                                               \tag{3}
\]
are monic of degree $M$ with $M$ distinct real roots in $(a,b)$.
Let $\mu_+$ and $\mu_-$ put weight one at each root of $P_+$ and
$P_-$ respectively. Then
\[
 m_k(\mu_+)=m_k(\mu_-)=m_k(\mu)
              \quad(0\leq k\leq M-q),                       \tag{4}
\]
where $\mu=\sum_h\delta_{p_h}$, and their Stieltjes transforms obey
\[
 \sigma_j\bigl(S_{\mu_+}(t_j)-S_{\mu_-}(t_j)\bigr)>0
                 \quad(1\leq j\leq q).                      \tag{5}
\]
The two root sets are not required to be disjoint from one another.
In particular $q\leq M/2$ permits $q$ arbitrary exterior signs while
preserving at least the first $q$ positive-degree moments and keeping
all atom weights equal to one.
\end{theorem}
\begin{proof}
Use (1) to choose $R$ with $(R/P)'(t_j)=\sigma_j$.
Choose disjoint closed intervals around the simple roots of $P$,
all inside $(a,b)$, with nonroot endpoints. At each interval the
endpoint values of $P$ have opposite signs. For all sufficiently
small $\delta>0$, $|\delta R|<|P|$ at every endpoint. Thus both
polynomials in (3) retain a sign change in every interval. The
intermediate value theorem gives $M$ distinct real roots; their
degrees show that these are all their roots, each simple.

Since $\deg R<q$, the coefficients of degrees $M-1,\ldots,q$ in
both polynomials agree with those of $P$. Newton's identities give
equality of their power sums of degrees $1,\ldots,M-q$. Degree
zero is the common number $M$ of unit-weight atoms. This proves (4).

The logarithmic derivative of a monic real-rooted polynomial is the
Stieltjes transform of its unit-weight root measure. Direct subtraction
therefore gives the exact identity
\[
 S_{\mu_+}(t)-S_{\mu_-}(t)
 =\frac{2\delta(R'P-RP')}{P^2-\delta^2R^2}
 =\frac{2\delta P^2}{P_+P_-}(R/P)'(t).                    \tag{6}
\]
At every exterior query, $P_+$ and $P_-$ are positive because all
their roots are below the query. The multiplier in the final
expression is strictly positive. Equation (5) follows exactly,
without replacing (6) by an asymptotic expansion.
\end{proof}

\paragraph{Quantitative load and the recursive gap.}
Applying a split with $q(M)$ queries to every node whose children
have $M$ poles would formally provide
$\sum_M (H/(2M))q(M)$ query slots. Taking $q(M)$ proportional to
$M/\log M$ would make this order $H\log\log H$. This is only a
slot count, not a construction: (3) replaces both children by
perturbations of the same prescribed polynomial $P$. It does not
preserve arbitrary separately prescribed child polynomials or their
descendant signs.

There is also a concrete coefficient-rigidity caution. If two child
polynomials have already been restricted to the forms
$Q_1^2-c_1$ and $Q_2^2-c_2$, with $Q_1,Q_2$ monic of degree $d$,
then agreement of their coefficients of degrees $2d-1$ down to $d$
forces $Q_1=Q_2$. Indeed, if the difference $Q_1-Q_2$ has degree
$e<d$ and is nonzero, their square difference has degree $d+e\geq d$.
Their difference is consequently constant after the stated agreement.
Thus starting from a hierarchy that already has only constant
square remainders can erase the apparent slack of the underloaded
local theorem. An actual recursion must keep additional polynomial
freedom or justify an adaptive partition of the poles.

\paragraph{Exact checks.}
\texttt{underloaded-equal-weight-split-check.py} uses rational poles,
queries, polynomials, and linear algebra. For $1\leq M\leq8$ and
every $1\leq q\leq M$, it checks (2) at a rational exterior point,
solves (1) for every sign word, chooses a rational $\delta$ by the
endpoint bound above, checks every endpoint sign and query sign,
and verifies the shared power sums through degree $M-q$ by Newton's
identities. No numerical root solver is used.
\end{document}
